\documentclass{proposal}

\title{NE511 Computational Project Notes}
\author{Kyle Hansen}
\date{\today}



\newcommand{\calL}{\mathcal{L}}
\newcommand{\calS}{\mathcal{S}}
\newcommand{\calH}{\mathcal{H}}
\newcommand{\calQ}{\mathcal{Q}}
\newcommand{\inv}{^{-1}}
\newcommand{\ddx}{\frac{\partial}{\partial x}}
\newcommand{\ddt}{\frac{\partial}{\partial t}}
\newcommand{\ra}{\rightarrow}
\newcommand{\qqquad}{\hskip 3em\relax}
\newcommand{\T}[1]{\text{{#1}}}
\newcommand{\p}{^\prime}
\newcommand{\rarr}{\rightarrow}

\usepackage{cleveref}
\usepackage[title]{appendix}




% \numberwithin{equation}{section}

\begin{document}
\maketitle

\tableofcontents

\section{Multiphysics Problem Description}

\subsection{Multigroup neutron transport equation with Backward Euler time discretization}

The problem described here is for 1-dimensional multigroup neutron transport, coupled to 1-dimensional advection-diffusion through temperature- and flux-dependent cross sections.

The one-dimensional time-dependent neutron transport equation (NTE) for prompt neutrons with isotropic scattering and isotropic source is

\begin{subequations}
    \begin{multline}
        \left[\frac{1}{v(E)}\ddt + \mu \ddx + \sigma_t(x, E) \right]\psi(x, \mu, E, t) =\\
        \frac{\chi}{2}\int\limits_{0}^{\infty}dE\p\; (1-\beta(E\p))\nu(x, E\p)\sigma_f(x, E\p) \phi(x, E\p, t) + \frac{1}{2}\int\limits_{0}^{\infty}dE\p\; \sigma_{s}(x, E\p \rarr E) \phi(x, E\p, t) + q(x, E, t), \\
         x \in [0, X], \; E \in [0,\infty), \;\mu \in [-1, 1], \; t>0
    \end{multline}
    \begin{equation}
        \phi(x, E, t) = \int\limits_{-1}^1 d\mu\; \psi(x, \mu, E, t)
    \end{equation}
    \begin{equation}
        \psi(x, \mu, E, 0) = \psi^0(x, \mu, E)
    \end{equation}
    \begin{equation}
        \psi(0, \mu, E, t) = \psi^{+}(\mu, E, t), \quad \mu > 0
    \end{equation}
    \begin{equation}
        \psi(X, \mu, E, t) = \psi^{-}(\mu, E, t), \quad \mu < 0.
    \end{equation}
\end{subequations}


Using the standard multigroup approximation, the equation for group $g$ becomes

\begin{subequations}
    \begin{multline}
        \left[\frac{1}{v_g}\ddt + \mu \ddx + \sigma_{t, g} \right]\psi_g(x, \mu, E, t) =\\
        \frac{\chi_g}{2}\sum_{g\p = 1}^{N_g} (1- \beta_{g\p})\nu\sigma_{f, g\p}(x) \phi_{g\p}(x,  t) + \frac{1}{2}\sum_{g\p = 1}^{N_g} \sigma_{s, g\p \rarr g}(x) \phi_{g\p}(x, t)+ q_{ g}(x,  t), \\
        g = 1, \dots, N_g,\; x \in [0, X],  \;\mu \in [-1, 1], \; t>0
    \end{multline}
    \begin{equation}
        \phi_g(x,  t) = \int\limits_{-1}^1 d\mu\; \psi_g(x, \mu, t)
    \end{equation}
    \begin{equation}
        \psi_g(x, \mu,  0) = \psi_g^0(x, \mu)
    \end{equation}
    \begin{equation}
        \psi_g(0, \mu,  t) = \psi_g^{+}(\mu,  t), \quad \mu > 0
    \end{equation}
    \begin{equation}
        \psi_g(X, \mu,  t) = \psi_g^{-}(\mu, t), \quad \mu < 0.
    \end{equation}
\end{subequations}

The Backward Euler temporal discretization scheme is

\begin{equation}
    \left.\ddt \psi(t)\right|_{t = t_{j+1}} \approx \frac{\psi(t_{j+1}) - \psi(t_j)}{\Delta t_{j}}
\end{equation}

where $t_{j=0} = 0$ and

\begin{equation}
    \Delta t_j = t_{j+1} - t_j.
\end{equation}

Substituting this into the multigroup neutron transport equation yields the effective steady-state problem at each time step,

\begin{subequations}\label{eq:equivalent-prompt-transport}
    \begin{multline}
        \left[ \mu \ddx + \hat{\sigma}^{(j+1)}_{t, g} \right]\psi^{(j+1)}_g(x, \mu, E) =\\
        \frac{\chi_g}{2}\sum_{g\p = 1}^{N_g} (1- \beta_{g\p})\nu\sigma^{(j+1)}_{f, g\p}(x) \phi^{(j+1)}_{g\p}(x) + \frac{1}{2}\sum_{g\p = 1}^{N_g} \sigma^{(j+1)}_{s, g\p \rarr g}(x) \phi_{g\p}^{(j+1)}(x)+ \hat{q}^{(j+1)}_{g}(x, \mu), \\
        g = 1, \dots, N_g,\; j=1, \dots N_j,\; x \in [0, X],  \;\mu \in [-1, 1]
    \end{multline}
    \begin{equation}
        \phi^{(j+1)}_g(x) = \int\limits_{-1}^1 d\mu\; \psi^{(j+1)}_g(x, \mu)
    \end{equation}
    \begin{equation}
        \psi^{(0)}_g(x, \mu) = \psi_g^0(x, \mu)
    \end{equation}
    \begin{equation}
        \psi^{(j+1)}_g(0, \mu) = \psi_g^{+}(\mu, t_{j+1}), \quad \mu > 0
    \end{equation}
    \begin{equation}
        \psi^{(j+1)}_g(X, \mu) = \psi_g^{-}(\mu, t_{j+1}), \quad \mu < 0
    \end{equation}
\end{subequations}

where the augmented total cross section and source are

\begin{equation}
    \hat{\sigma}^{(j+1)}_{t, g} = \left(\sigma^{(j+1)}_{t, g}(x) +\frac{1}{v_g \Delta t_j} \right)
\end{equation}

\begin{equation}
    \hat{q}_g^{(j+1)}(x, \mu) = \left(q_g(x, t_{j+1}) + \frac{1}{v_g \Delta t_j}\psi_g^{(j)}(x, \mu) \right).
\end{equation}

\subsection{Precursor balance equations}

In one dimension, the precursor balance equation is 

\begin{subequations}
    \begin{equation}
        \ddt C_{d, k}(x, t) = -\lambda_k C_{d, k}(x, t) + \sum_{g = 1}^{N_g}\beta_{k, g}(x, t)\nu\sigma_{f, g}(x)\phi_g(x, t), \qquad k= 1,\dots, 6, \; t>0
    \end{equation}
    \begin{equation}
        C_{d, k}(x, 0) = C_{d, k}^{0}(x).
    \end{equation}
\end{subequations}

With the Backward Euler temporal discretization, this may be written, after some algebra, as

\begin{equation}\label{eq:precursor-algebra}
    C_{d, k}^{(j+1)}(x) = \frac{\frac{C_{d, k}^{(j)}}{\Delta t_j} + \sum_{g=1}^{N_g} \beta_{g, k}\nu\sigma_{f, g}\phi_g^{(j+1)}}{\lambda_k + \frac{1}{\Delta t_j}}.
\end{equation}

In the neutron transport equation, the delayed neutron source is

\begin{equation}\label{eq:delayed-source-definition}
    q_{del, g}(x, \mu) = \frac{1}{2}\sum_{k=1}^{6}\lambda_k \chi_{g, k}^d C_{d, k},
\end{equation}

so, substituting \cref{eq:precursor-algebra} into \cref{eq:delayed-source-definition}, the time-discretized delayed neutron source in the transport equation becomes

\begin{equation}\label{eq:delayed-source}
    q_{del, g}(x, \mu) = \frac{1}{2}\sum_{g\p = 1}^{N_g} {\Gamma}^{(j+1)}_{g\p\rarr g} \nu \sigma_{f, g\p} \phi_{g\p}^{(j+1)} + \overline{q}^{(j+1)}_{del, g}(x, \mu)
\end{equation}

where

\begin{equation}
    {\Gamma}^{(j+1)}_{g\p\rarr g} = \sum_{k=1}^6 \chi_{g, k}\frac{\lambda_k\Delta t_j \beta_{g\p, k}}{\lambda_k \Delta t_j + 1}
\end{equation}

\begin{equation}
    \overline{q}^{(j+1)}_{del, g}(x, \mu) = \frac{1}{2}\sum_{k=1}^6 \lambda_k \chi_{g, k} \left(\frac{C^{(j)}}{\lambda_k \Delta t_j + 1} \right).
\end{equation}

Here, $\overline{q}^{(j+1)}$ represents the decay of precursors that existed in time step $j$, and ${\Gamma}^{(j+1)}_{g\p\rarr g}$ represents the proportion of fissions in group $g\p$, during the time step $\Delta t_j$, whose delayed neutrons decay during the same time step (immediate precursor decay).

\subsection{Equivalent steady-state transport problem with delayed neutrons}

Adding \cref{eq:delayed-source} to \cref{eq:equivalent-prompt-transport,}, the quasi-steady state transport equation to be solved at each time step, including the effect of delayed neutrons, is

\begin{subequations}\label{eq:equivalent-steady-with-delayed}
   \begin{multline}
        \left[ \mu \ddx + \hat{\sigma}^{(j+1)}_{t, g} \right]\psi^{(j+1)}_g(x, \mu, E) =\\
        \frac{\chi_g}{2}\sum_{g\p = 1}^{N_g} (1- \beta_{g\p})\nu\sigma^{(j+1)}_{f, g\p}(x) \phi^{(j+1)}_{g\p}(x) + \frac{1}{2}\sum_{g\p = 1}^{N_g} \left(\sigma^{(j+1)}_{s, g\p \rarr g}(x) + {\Gamma}^{(j+1)}_{g\p\rarr g} \nu \sigma_{f, g\p} \right)\phi_{g\p}^{(j+1)}(x)\\
          +\hat{q}^{(j+1)}_{g}(x, \mu)+ \overline{q}^{(j+1)}_{del, g}(x, \mu), \\
        g = 1, \dots, N_g,\; j=1, \dots N_j,\; x \in [0, X],  \;\mu \in [-1, 1]
    \end{multline}
    \begin{equation}
        \phi^{(j+1)}_g(x) = \int\limits_{-1}^1 d\mu\; \psi^{(j+1)}_g(x, \mu)
    \end{equation}
    \begin{equation}
        \psi^{(j+1)}_g(0, \mu) = \psi_g^{+}(\mu, t_{j+1}), \quad \mu > 0
    \end{equation}
    \begin{equation}
        \psi^{(j+1)}_g(X, \mu) = \psi_g^{-}(\mu, t_{j+1}), \quad \mu < 0
    \end{equation}
\end{subequations}

where

\begin{subequations}
    \begin{equation}
        \psi^{(0)}_g(x, \mu) = \psi_g^0(x, \mu), \qquad g= 1, \dots, N_g
    \end{equation}
    \begin{eqnarray}
        C_{d, k}^{(0)}(x) = C_{d, k}^0(x), \qquad k = 1,\dots, 6
    \end{eqnarray}
\end{subequations}



After solving for the neutron distribution, the precursor distribution can be recovered using \cref{eq:precursor-algebra}.


\subsection{Multigroup Low-Order Quasidiffusion Equations}

Taking the zeroth and first spatial moments ($\int \cdot \,d\mu$ and $\int \mu \cdot d\mu$) of \cref{eq:equivalent-steady-with-delayed} (omitting $j+1$ index) yields

\begin{subequations}
    \begin{multline}
        \ddx J_g(x) + \sigma_{t, g}\phi_g(x) = \chi_g \overline{(1-\beta)\nu\sigma_{f}}\phi \\
        +  \sum_{g\p = 1}^{N_g} \left(\sigma_{s, g\p \rarr g} + {\Gamma}_{g\p\rarr g} \nu \sigma_{f, g\p} \right)\phi_{g\p}(x) + \hat{Q}_g(x) + \overline{Q}_{del, g}
    \end{multline}
    \begin{equation}
        \ddx (E_g \phi_g) + \sigma_{t, g}J_g = \tilde{q}_g
    \end{equation}
    \begin{equation}
        J_g(0) = C^b_{0, g}\phi_g(0), \qquad J_g(X) = C^b_{X, g}\phi_g(X)
    \end{equation}
    
\end{subequations}

where the initial conditions are derived from taking the zeroth and first spatial moments of $\psi^0$ and

\begin{equation}
    \hat{Q}_g(x) = \left(2 q_g(x) + \frac{1}{v_g \Delta t_j}\phi_g^{(j)}(x) \right)
\end{equation}
\begin{equation}
    \overline{Q}_{del, g} = \sum_{k=1}^6 \lambda_k \chi_{g, k} \left(\frac{C^{(j)}}{\lambda_k \Delta t_j + 1} \right)
\end{equation}
\begin{equation}
    \tilde{q}_g = \frac{1}{v_g \Delta t_j}J^{(j)}(x)
\end{equation}
        
and

\begin{equation}
    \overline{(1-\beta)\nu\sigma_{f}} = \frac{\sum_{g=1}^{N_g}(1-\beta_g)\nu\sigma_{f, g}}{\sum_{g=1}^{N_g} \phi_g}
\end{equation}

\begin{equation}
    \phi = \sum_{g=1}^{N_g} \phi_g.
\end{equation}

The nonlinear closure term (Eddington Factor) is

\begin{equation}
    E_g = \frac{\int\limits_{-1}^1 \mu^2 \psi_g d \mu}{\int\limits_{-1}^1 \psi_g d\mu}
\end{equation}

\subsection{Advection-diffusion equation}

The advection-diffusion equation in 1D is


\begin{equation}\label{eq:timedep-adv-diffusion}
    \left[ \ddt - \ddx k_b \ddx + \dot{m}c_p \ddx \right]T(x, t) = W_f f(x, t)
\end{equation}

or in operator form,

\begin{equation}
    \calH T(x, t) = Q_f
\end{equation}

with appropriate boundary conditions. the fission reaction rate $f$ is 

\begin{equation}
    f(x) = \sum_{g=1}^{N_g} \sigma_{f, g}\phi_g(x, t)
\end{equation}

and

\begin{description}[labelindent=1cm, labelwidth=1.2cm, before={\renewcommand\makelabel[1]{##1 $=$}}]
    \item[$c_p$] specific heat capacity
    \item[$k_b$] turbulent diffusion coefficient
    \item[$\dot{m}$] mass flow rate
    \item[$W_f$] fission energy deposition
\end{description}

Discretizing with Backward Euler, this equation becomes

\begin{align}
    \left[ \frac{1}{\Delta t_j} - \ddx k_b \ddx + \dot{m}c_p \ddx \right]T^{(j+1)}(x) &= W_f f^{(j+1)}(x) + \frac{1}{\Delta t_j}T^{(j)}\\
    \hat{\calH}T^{(j+1)}(x) &= \hat{Q}_f^{(j+1)}
\end{align}

which is an equivalent steady-state advection-diffusion problem at each time step with a local removal term.

\subsection{Cross-section feedback}

For this study, the macroscopic feedback will be modeled with

\begin{equation}\label{eq:temp-xs}
    \sigma_u(x) = \sigma_{u, 0} + \sigma_{u, 1, d}\left( \sqrt{T(x) + \kappa W_f f(x)} - \sqrt{T_{d, 0}} \right)  + \sigma_{u, 1, b}\left(T(x) -T_0 \right)
\end{equation}

for each reaction $u$, where  $\kappa, \sigma_{u, 0}, \sigma_{u, 1, d}, \sigma_{u, 1, b}, T_{d, 0},$ and $T_0$ are material constants, defined for each neutron energy group $g$.

\section{MNP Method}

\subsection{Basic method details}

\subsection{Iterative Algorithms}


\section{NILO-MNP method}

\subsection{Method details}

\subsection{Iterative Algorithm}

\clearpage
\begin{appendices}

\section{LD Discretization of Neutron Transport equation}

\section{FV Discretization of advection-diffusion equation}

\section{LD Discretization of QD equations}

\end{appendices}










\end{document}